\documentclass[11pt,a4paper]{article}

\usepackage[utf8]{inputenc}
\usepackage[margin=2.2cm]{geometry}
\usepackage{amsmath,amssymb}
\usepackage{booktabs}
\usepackage{tabularx}
\usepackage{xcolor}
\usepackage{listings}
\usepackage{hyperref}
\usepackage{microtype}
\usepackage{enumitem}

\definecolor{codegreen}{rgb}{0,0.6,0}
\definecolor{codegray}{rgb}{0.5,0.5,0.5}
\definecolor{codepurple}{rgb}{0.58,0,0.82}
\definecolor{backcolour}{rgb}{0.97,0.97,0.98}
\definecolor{darkblue}{rgb}{0.0,0.2,0.6}

\lstset{
    backgroundcolor=\color{backcolour},
    commentstyle=\color{codegreen},
    keywordstyle=\color{darkblue}\bfseries,
    numberstyle=\tiny\color{codegray},
    stringstyle=\color{codepurple},
    basicstyle=\ttfamily\footnotesize,
    breakatwhitespace=false,
    breaklines=true,
    captionpos=b,
    keepspaces=true,
    numbers=left,
    numbersep=5pt,
    showspaces=false,
    showstringspaces=false,
    showtabs=false,
    tabsize=2,
    frame=single,
    rulecolor=\color{lightgray}
}

\hypersetup{
    colorlinks=true,
    linkcolor=darkblue,
    urlcolor=darkblue,
    citecolor=darkblue,
    pdftitle={Laporan Revisi & Desain Ulang PhysLang v2},
    pdfauthor={Anaf & Xinxinran}
}

\title{\textbf{\LARGE Laporan Komprehensif Revisi \& Desain Ulang PhysLang v2}\\[0.3em]
\Large Transformasi Menuju Domain-Specific Language (DSL) Fisika Komputasional}
\author{\textbf{Anaf (Lead Architect)} \and \textbf{Xinxinran (Co-Compiler Analyst)}}
\date{November 2026}

\begin{document}

\maketitle
\begin{abstract}
Dokumen ini menyajikan analisis menyeluruh, justifikasi teoritis, tata bahasa formal, serta rancangan teknis transisi \textbf{PhysLang} dari implementasi awal (\textit{v1}) menuju \textbf{PhysLang v2}. Desain ulang ini dilakukan menyusul evaluasi akademik yang mengidentifikasi bahwa tata bahasa permukaan PhysLang v1 terlalu menyerupai bahasa pemrograman imperatif umum (\textit{general-purpose}) dan belum menonjolkan abstraksi domain fisis secara kuat. PhysLang v2 mengadopsi paradigma \textit{Computational Physicist POV} (gaya lembar kerja laboratorium) dengan menghilangkan \textit{boilerplate} sintaksis yang tidak perlu, memperkenalkan pendefinisian rumus fisis sebagai warga kelas satu (\textit{first-class formula declaration} \texttt{form}), proyeksi konversi satuan (\texttt{in}), serta percabangan kondisional (\texttt{if-else}) yang divalidasi ketat oleh matriks dimensi vektor $\mathbb{Z}^7$. Seluruh perubahan tata bahasa, aturan produksi (EBNF), struktur kuadrupel, tabel perbandingan, dan pemetaan perubahan modul kode sumber disajikan secara mendalam.
\end{abstract}

\tableofcontents
\newpage

\section{Latar Belakang dan Dekonstruksi Masukan Evaluator}

\subsection{Kritik Akademik terhadap PhysLang v1}
Pada tahap pengujian awal tata bahasa, evaluator memberikan catatan kritis yang mencakup:
\begin{enumerate}[noitemsep]
    \item Sintaksis yang diajukan terkesan seperti bahasa pemrograman umum (\textit{general-purpose language}) seperti JavaScript atau Kotlin mini.
    \item Kurang menonjolnya abstraksi domain fisika pada tingkat aturan produksi tata bahasa permukaan (\textit{surface grammar}).
    \item Ketiadaan konstruksi fungsional atau formulasi eksplisit seperti pendefinisian entitas $\text{Area}(a, b)$ yang lumrah ditemukan pada bahasa berbasis domain.
    \item Ketiadaan struktur alur kendali (\textit{control flow}) yang membuat bahasa terlihat belum memiliki kapabilitas pemrosesan bahasa terstruktur secara komprehensif.
\end{enumerate}

\subsection{Analisis Akar Masalah: Disonansi Sintaksis Permukaan vs Semantik Internal}
Secara teknis, \textit{backend} PhysLang v1 yang dibangun menggunakan Kotlin sebenarnya telah memiliki analisis semantik fisis yang sangat maju, di antaranya representasi aljabar linier vektor dimensi SI 7-elemen ($\mathbb{Z}^7$) serta validasi kompatibilitas dimensi pada pohon sintaks (\textit{AST}).

Namun, seorang penilai tata bahasa kompilator mengevaluasi bahasa melalui \textbf{aturan sintaksis (EBNF)} yang tercantum pada dokumen laporan. Pada PhysLang v1, aturan produksi yang tampak adalah:
\begin{align*}
    \text{VarDeclStmt} &\to \texttt{"var"}\ \text{IDENTIFIER}\ ( \texttt{"="}\ \text{Expression} )^? \\
    \text{PrintStmt} &\to \texttt{"print"}\ \text{Expression}
\end{align*}
Keberadaan kata kunci \texttt{"var"} dan \texttt{"print"} secara visual menyamarkan hakikat domain fisika menjadi sekadar program imperatif biasa yang kebetulan memiliki satuan pada literal angka.

\subsection{Studi Komparasi Bahasa Domain: Dari R hingga Wolfram dan Frink}
Dalam perancangan DSL, struktur bahasa harus mencerminkan cara berpikir praktisi domain tersebut:
\begin{itemize}[noitemsep]
    \item \textbf{Bahasa Statistik (R):} Memiliki operator formula khusus \texttt{\textasciitilde} (\textit{Wilkinson-Rogers notation}) untuk merepresentasikan model regresi linear $\text{mpg} \sim \text{wt} + \text{cyl}$. Karakteristik statistik dijadikan bagian dari tata bahasa, bukan sekadar pemanggilan fungsi pustaka.
    \item \textbf{Modelica:} Standar industri pemodelan fisika berbasis relasi kesetimbangan persamaan (\textit{equation-oriented}), di mana variabel mengikat dimensi fisis tertentu (\texttt{Length}, \texttt{Force}) dan mengeliminasi konsep mutasi variabel imperatif.
    \item \textbf{Wolfram Language \& Frink:} Memperlakukan konversi satuan dan operasi dimensi sebagai operasi tingkat pertama (\textit{first-class operator}) menggunakan sintaks natural seperti \texttt{x in meter} atau operator proyeksi fisis.
\end{itemize}

PhysLang v2 mengambil posisi sebagai \textbf{Lightweight Interactive Physics DSL} yang mengedepankan kemudahan penulisan formula di laboratorium tanpa beban \textit{solver} numerik yang rumit.

\section{Filosofi Desain PhysLang v2: Sudut Pandang Fisikawan}

PhysLang v2 membuang paradigma pemrogram perangkat lunak (\textit{software engineer}) yang kaku dan menggantinya dengan paradigma fisikawan laboratorium:
\begin{enumerate}
    \item \textbf{Deklarasi Variabel Nol-Boilerplate (\texttt{id = expr}):}
    Fisikawan di papan tulis tidak pernah menulis \texttt{var m = 10 kg}. Mereka menulis $m = 10\text{ kg}$. Oleh karena itu, kata kunci \texttt{var} dihapus total dan digantikan oleh penugasan langsung.
    
    \item \textbf{Deklarasi Rumus Fisis Eksplisit (\texttt{form}):}
    Untuk mendefinisikan rumus seperti $\text{Area}(a, b)$, fisikawan tidak memerlukan kurung kurawal berlapis atau kata kunci \texttt{return}. PhysLang v2 memperkenalkan konstruksi \texttt{form}:
    \begin{equation*}
        \texttt{form Area(p, l) = p * l}
    \end{equation*}
    Penggunaan kata kunci \texttt{form} menghilangkan seluruh potensi ambiguitas pada parser antara pemanggilan fungsi (\textit{CallExpr}) dan deklarasi rumus baru.
    
    \item \textbf{Pelestarian Fitur Kustom Satuan (\texttt{unit}):}
    Fitur unggulan pembuatan satuan dinamis berbasis ekspresi fisis tetap dipertahankan dan disempurnakan:
    \begin{equation*}
        \texttt{unit kph = 1000 m / 3600 s}
    \end{equation*}
    
    \item \textbf{Operator Proyeksi Satuan (\texttt{in}):}
    Hasil evaluasi dapat langsung diamati dalam skala satuan tertentu secara deklaratif:
    \begin{equation*}
        \texttt{inspect Ek in kJ}
    \end{equation*}
    
    \item \textbf{Kondisi Batas Fisis (\texttt{if-else}):}
    Struktur alur kendali dihadirkan bukan untuk rekayasa algoritma umum, melainkan untuk mengevaluasi kondisi batas fisis (\textit{physical boundary conditions}), titik kritis suhu, atau transisi fase, di mana operator pembanding ($<, >, ==$) wajib berdimensi setara.
\end{enumerate}

\section{Tata Bahasa Formal (EBNF Grammar)}

Berikut adalah spesifikasi tata bahasa bebas konteks (\textit{Context-Free Grammar}) dalam notasi \textit{Extended Backus-Naur Form} (EBNF) untuk PhysLang v2.

\subsection{Aturan Leksikal (Tokens)}
\begin{lstlisting}[language=SQL]
DIGIT            ::= [0-9]
INT_LITERAL      ::= DIGIT+
DOUBLE_LITERAL   ::= DIGIT+ "." DIGIT+
NUMBER_LITERAL   ::= INT_LITERAL | DOUBLE_LITERAL
STRING_LITERAL   ::= '"' [^"]* '"'
IDENTIFIER       ::= [a-zA-Z_][a-zA-Z0-9_]*

KEYWORD          ::= "unit" | "form" | "inspect" | "in" | "if" | "else" | "null"

ASSIGN_OP        ::= "=" | "+=" | "-=" | "*=" | "/="
COMPARE_OP       ::= "==" | "!=" | "<" | "<=" | ">" | ">="
ARITH_OP         ::= "+" | "-" | "*" | "/" | "^"
DELIMITER        ::= "(" | ")" | "{" | "}" | "," | "\n"
\end{lstlisting}

\subsection{Aturan Sintaksis (Syntactic Production Rules)}
\begin{lstlisting}[language=SQL]
Program          ::= ( Statement | "\n" )* EOF

Statement        ::= UnitDeclStmt
                   | FormulaDeclStmt
                   | IfStmt
                   | InspectStmt
                   | BlockStmt
                   | VarAssignStmt
                   | ExprStmt

UnitDeclStmt     ::= "unit" IDENTIFIER "=" Expression
FormulaDeclStmt  ::= "form" IDENTIFIER "(" ParameterList? ")" "=" Expression
ParameterList    ::= IDENTIFIER ( "," IDENTIFIER )*

IfStmt           ::= "if" "(" ConditionExpr ")" BlockStmt ( "else" BlockStmt )?
ConditionExpr    ::= Expression COMPARE_OP Expression

InspectStmt      ::= "inspect" Expression ( "in" IDENTIFIER )?
BlockStmt        ::= "{" ( Statement | "\n" )* "}"
VarAssignStmt    ::= IDENTIFIER ASSIGN_OP Expression
ExprStmt         ::= Expression

Expression       ::= ConversionExpr
ConversionExpr   ::= TermAddSub ( "in" IDENTIFIER )?
TermAddSub       ::= TermMulDiv ( ( "+" | "-" ) TermMulDiv )*
TermMulDiv       ::= TermPower ( ( "*" | "/" ) TermPower )*
TermPower        ::= UnaryExpr ( "^" UnaryExpr )*
UnaryExpr        ::= "-" UnaryExpr | PostfixExpr
PostfixExpr      ::= Primary ( "(" ArgumentList? ")" )*
ArgumentList     ::= Expression ( "," Expression )*

Primary          ::= QuantityLiteral
                   | NUMBER_LITERAL
                   | STRING_LITERAL
                   | IDENTIFIER
                   | "(" Expression ")"

QuantityLiteral  ::= NUMBER_LITERAL IDENTIFIER
\end{lstlisting}

\section{Struktur Kuadrupel Formal: $G = (V_N, V_T, S, P)$}

Berdasarkan definisi matematis Chomsky, tata bahasa PhysLang v2 didefinisikan oleh kuadrupel:
\begin{equation}
    G = (V_N, V_T, S, P)
\end{equation}

\subsection{1. Himpunan Simbol Non-Terminal ($V_N$)}
Entitas variabel tata bahasa yang digunakan dalam penurunan struktur sintaksis:
\begin{align*}
V_N = \{ & \text{Program}, \text{Statement}, \text{UnitDeclStmt}, \text{FormulaDeclStmt}, \text{ParameterList}, \\
         & \text{IfStmt}, \text{ConditionExpr}, \text{InspectStmt}, \text{BlockStmt}, \text{VarAssignStmt}, \\
         & \text{ExprStmt}, \text{Expression}, \text{ConversionExpr}, \text{TermAddSub}, \text{TermMulDiv}, \\
         & \text{TermPower}, \text{UnaryExpr}, \text{PostfixExpr}, \text{ArgumentList}, \text{Primary}, \\
         & \text{QuantityLiteral} \}
\end{align*}

\subsection{2. Himpunan Simbol Terminal ($V_T$)}
Token leksikal fisik yang dikenali langsung oleh penganalisis leksikal:
\begin{align*}
V_T = \{ & \text{NUMBER\_LITERAL}, \text{STRING\_LITERAL}, \text{IDENTIFIER}, \\
         & \texttt{"unit"}, \texttt{"form"}, \texttt{"inspect"}, \texttt{"in"}, \texttt{"if"}, \texttt{"else"}, \texttt{"null"}, \\
         & \texttt{"="}, \texttt{"+="}, \texttt{"-="}, \texttt{"*="}, \texttt{"/="}, \\
         & \texttt{"=="}, \texttt{"!="}, \texttt{"<"}, \texttt{"<="}, \texttt{">"}, \texttt{">="}, \\
         & \texttt{"+"}, \texttt{"-"}, \texttt{"*"}, \texttt{"/"}, \texttt{"\string^"}, \\
         & \texttt{"("}, \texttt{")"}, \texttt{"\{"}, \texttt{"\}"}, \texttt{","}, \texttt{"\textbackslash n"} \}
\end{align*}

\subsection{3. Simbol Awal ($S$)}
\begin{equation*}
    S = \text{Program}
\end{equation*}

\subsection{4. Himpunan Aturan Produksi ($P$)}
Himpunan relasi penurunan produksi yang memetakan elemen $V_N \to (V_N \cup V_T)^*$, persis sebagaimana didefinisikan pada subbab 3.2.

\section{Contoh Komparasi Kode Program: v1 vs v2}

Berikut adalah perbandingan konkret skrip yang mengevaluasi dinamika gerak dan gaya:

\subsection{Implementasi PhysLang v1 (Imperatif Konvensional)}
\begin{lstlisting}
unit kph = 1000 m / 3600 s

var m = 1200 kg
var v = 90 kph
var a = 2.5 m/s^2

var F = m * a
var Ek = 0.5 * m * v^2

print "Gaya resultan: " + F
print "Energi kinetik: " + Ek
\end{lstlisting}

\subsection{Implementasi PhysLang v2 (Sintaks Laboratorium Fisika Murni)}
\begin{lstlisting}
# 1. Deklarasi satuan pengukuran eksperimen
unit kph = 1000 m / 3600 s
unit Newton = 1 kg * m / s^2

# 2. Perumusan hukum fisika dasar (First-class Formula)
form Area(p, l) = p * l
form Gaya(m, a) = m * a
form EnergiKinetik(m, v) = 0.5 * m * v^2

# 3. Lembar data observasi (tanpa var)
massaUji = 1200 kg
kecepatan = 90 kph
percepatan = 2.5 m/s^2

# 4. Komputasi menggunakan formula
F = Gaya(massaUji, percepatan)
Ek = EnergiKinetik(massaUji, kecepatan)

# 5. Observasi dan proyeksi satuan
inspect F in Newton
inspect Ek in kJ

# 6. Pengujian kondisi batas sistem mekanik
batasKritis = 2500 Newton
if (F > batasKritis) {
    inspect "Peringatan: Terjadi over-stress mekanis!"
} else {
    inspect "Status beban normal."
}
\end{lstlisting}

\section{Daftar Fitur: Penambahan, Modifikasi, dan Penghapusan}

\begin{table}[h]
\centering
\small
\begin{tabularx}{\textwidth}{lp{4.2cm}X}
\toprule
\textbf{Nama Fitur} & \textbf{Klasifikasi} & \textbf{Deskripsi Teknis \& Rasionalisasi} \\
\midrule
\texttt{form} & \textbf{Fitur Baru} & Mendefinisikan fungsi fisis satu baris matematis murni. Menghilangkan ambiguitas pemanggilan fungsi dengan deklarasi formula. \\
\texttt{in} & \textbf{Fitur Baru} & Operator konversi satuan langsung. Memproyeksikan besaran $\mathbb{Z}^7$ ke unit target yang kompatibel dimensinya. \\
\texttt{if-else} & \textbf{Fitur Baru} & Kontrol alur untuk evaluasi kondisi batas fisis. Memeriksa relasi komparasi dengan verifikasi dimensi wajib setara. \\
\texttt{inspect} & \textbf{Modifikasi} & Menggantikan kata kunci \texttt{print} agar merefleksikan terminologi observasi sains di laboratorium. \\
\texttt{id = expr} & \textbf{Modifikasi} & Mengizinkan inisialisasi variabel langsung tanpa kata kunci pembuka deklarasi. \\
\texttt{var} & \textbf{Dihapus} & Dieliminasi total karena merupakan artefak bahasa pemrograman umum yang merusak estetika DSL sains. \\
\texttt{unit} & \textbf{Dipertahankan \& Diperkuat} & Fitur inti perhitungan aljabar satuan kustom tetap dipertahankan 100\% tanpa perubahan mekanisme komputasi. \\
\bottomrule
\end{tabularx}
\caption{Matriks Perubahan Fitur Bahasa PhysLang v2}
\end{table}

\section{Pemetaan Perubahan pada Arsitektur Kompilator Kotlin}

Implementasi PhysLang v2 tetap mempertahankan arsitektur \textit{Pratt Parser} dan \textit{Tree-Walking Interpreter} yang ada, namun memperbarui beberapa modul kunci:

\subsection{1. Modul Leksikal (\texttt{lang.lexer})}
\begin{itemize}[noitemsep]
    \item \textbf{Penambahan Token Type:} \texttt{Form}, \texttt{Inspect}, \texttt{In}, \texttt{If}, \texttt{Else}, \texttt{Less}, \texttt{Greater}, \texttt{LessEqual}, \texttt{GreaterEqual}, \texttt{EqualEqual}, \texttt{NotEqual}.
    \item \textbf{Penghapusan Token Type:} Menandai \texttt{Var} sebagai \textit{deprecated} atau dihapus dari daftar kata kunci di \texttt{Keywords.kt}.
\end{itemize}

\subsection{2. Modul AST (\texttt{lang.ast})}
\begin{itemize}[noitemsep]
    \item \textbf{Penambahan Node Statement:}
    \begin{itemize}
        \item \texttt{FormulaDeclStmt(name: Token, params: List<Token>, body: Expr)}
        \item \texttt{IfStmt(condition: Expr, thenBranch: Stmt, elseBranch: Stmt?)}
        \item \texttt{InspectStmt(expr: Expr, targetUnit: Token?)}
    \end{itemize}
    \item \textbf{Penyesuaian Statement Penugasan:} Menyatukan deklarasi dan penugasan ke dalam \texttt{AssignStmt} yang secara otomatis mendaftarkan simbol ke \texttt{Environment} jika belum ada.
\end{itemize}

\subsection{3. Modul Parser (\texttt{lang.parser})}
\begin{itemize}[noitemsep]
    \item Mengimplementasikan metode \texttt{parseFormulaDecl()} untuk mengenali pola \texttt{form IDENTIFIER ( args ) = expr}.
    \item Mengimplementasikan metode \texttt{parseIfStmt()} yang mem-parse kondisi dalam kurung dan blok eksekusi cabang.
    \item Menambahkan operator infix \texttt{"in"} ke dalam tabel presedensi (\textit{binding power}) ekspresi.
\end{itemize}

\subsection{4. Modul Penganalisis Semantik (\texttt{lang.semantic.SemanticAnalyzer})}
\begin{itemize}[noitemsep]
    \item \textbf{Verifikasi Komparasi:} Memastikan bahwa ekspresi di sisi kiri dan kanan operator perbandingan ($<, >, ==$) memiliki vektor eksponen dimensi $\mathbb{Z}^7$ yang identik. Membandingkan panjang dengan massa akan langsung memicu \texttt{DiagnosticType.ERROR}.
    \item \textbf{Verifikasi Proyeksi \texttt{in}:} Memeriksa apakah satuan target pada operasi \texttt{in} memiliki dimensi yang kompatibel dengan ekspresi sumber.
\end{itemize}

\subsection{5. Modul Runtime Evaluator (\texttt{lang.ast.visitors.PhysicsEvaluator})}
\begin{itemize}[noitemsep]
    \item Menyimpan objek \texttt{FormulaValue} di \texttt{Environment}. Ketika formula dipanggil via \texttt{CallExpr}, dibuat \textit{local environment} untuk mengevaluasi parameter aktual dan mengembalikan \texttt{PhysicsValue}.
    \item Mengevaluasi cabang \texttt{IfStmt} berdasarkan nilai kebenaran ekspresi boolean.
    \item Menghitung faktor skala konversi satuan saat mengevaluasi pernyataan \texttt{InspectStmt} dengan klausa \texttt{in}.
\end{itemize}

\section{Kesimpulan}
Transformasi PhysLang menuju versi 2 berhasil merekonsiliasi ekspektasi akademik evaluator kompilator dengan fondasi komputasi yang telah dibangun. Dengan mengadopsi \textit{Domain-Specific Grammar} yang dirancang khusus untuk fisikawan, PhysLang v2 memiliki identitas yang kuat, bebas dari kesan bahasa umum tiruan, namun tetap mempertahankan keunggulan perhitungan dimensi otomatis dan fleksibilitas pendefinisian satuan mandiri.

\end{document}
